\section{Setting, protocol, and evidence}
\label{sec:methods}

\subsection{Task and problem source}

The full user objective is reproduced in Appendix~\ref{app:protocol}.
It specifies an overall time limit and a research objective, rather
than a list of preselected questions or a target number of outputs.
The agent's initial plan explicitly prioritized precise statements,
literature checks, finite counterexamples, elementary reductions, and
independent implementations of computational checks. It also required
partial results and unsuccessful searches to be recorded.

The input file \artifact{docs/21tkt.pdf} has SHA-256 digest
\begin{center}
\small\path{2fcce9b98a4df10267fe120229217bfe556c70510704e311540da0cef438f911}.
\end{center}
The agent extracted a main-body index with 1,308 entries, including
all 150 entries of Issue 21. This is an index of problem statements,
not a verified count of open problems. Editorial stars, solved
subparts, duplicates, and literature published after a problem was
posed all affect the open status of individual assertions.
The source PDF, rather than its text extraction, was treated as
authoritative for mathematical notation and scope.

\subsection{Agent and computational environment}

The experiment used Codex 0.153.4 with \texttt{gpt-6-astra} and
reasoning effort \texttt{xhigh}. The full JSONL record corroborates the
settings supplied during manuscript preparation: all 139 retained
turn-context records identify this model and effort, including 133 in
the research window. The session metadata records the OpenAI provider
and the CLI application version. These observations identify the
recorded configuration, without establishing every provider-side
implementation detail.

The agent worked in a persistent Linux workspace with shell access,
file editing, Git, internet source retrieval, and tools for inspecting
PDF pages. GAP 4.16.1 supplied the main finite-group computations.
The retained scripts also use Python and C/C++ for exact arithmetic,
word calculations, independent finite models, enumeration, and
certificate parsing. The source and report files record which
libraries or catalogues each computation depends on. Versions not
retained in the experimental records cannot be reconstructed merely
by inspecting the later manuscript environment.

The user initially supplied an inaccessible GAP path and repaired it
during setup. The visible human message is \emph{sorry, gap is now here}.
The subsequent reports establish that GAP was available with the
required packages. The initial suggestion that a replacement
installation might be needed is therefore not evidence that such an
installation actually occurred.

The working limits were generally 20 CPU cores and 100\,GB of memory.
The repository contains bounded launchers and process observations,
but not a continuous, comprehensive record of global peak resource
use. In the largest retained certificate check, six GAP workers each
had a 14\,GiB workspace ceiling, for a combined configured ceiling of
84\,GiB. A workspace ceiling is not an observed resident-memory peak.
Likewise, local CPU limits do not measure the remote compute used by
the language model.

The research plan prohibited delegation to other agents without
authorization, and the reports record no delegated agents. Parallel
GAP or C++ processes were computational workers, not separate language
model researchers. The same agent developed the arguments, wrote
checkers, and conducted the internal reviews. The phrase
\emph{independent implementation} below therefore refers to a separate
representation or algorithm, not to an independent human author or
an independently trained reviewing model.

\subsection{Persistence and local records}

The workspace served as external memory. Plans and status files recorded
the next steps; proof drafts preserved mathematical arguments; scripts
and logs preserved computations; and Git recorded successive states.
The research window contains 54 task-start events, 52 task-completion
events, one aborted turn, and 79 context compactions. One further task
completion occurs during closeout. The opening task was paused during
setup; after the GAP correction at 20:58:51 UTC it resumed at
20:59:58 UTC. The wall-clock budget includes this interval. Repeated
continuations and compactions are visible in the record; the work was
not a single uninterrupted model response.

There were 198 commits during the 48-hour window. The first, at
21:00:37 UTC on 10 September, initialized the investigation.
The 199th commit, \snapshot{}, completed the local handoff at
20:57:02 UTC on 12 September. Commit times date repository checkpoints;
they do not directly date the discovery of an idea or establish when
its full final scope was proved.

At \snapshot{}, the Git tree contains 4,147 tracked files with a total
logical size of 2,440,677,642 bytes. This includes 430 files under
\artifact{research/}, 456 under \artifact{scripts/}, and 2,755 under
\artifact{results/}, as well as references, the source Notebook, and
other records. Logical file size is not compressed repository size.
The manuscript analysis reads this fixed revision, so later writing
and review do not increase the experimental counts.

\subsection{Session records and observable operations}

The full JSONL session supplied on 15 September contains 36,016
records, from 10 September 20:56:46.325 UTC through 14 September
17:05:09.077 UTC. It includes the research run, its closeout, and
initial manuscript preparation. We separate these phases before
counting operations or usage. Manuscript preparation begins with the
new goal recorded at 14 September 14:12:56.892 UTC.

The earlier \artifact{session.html} and \artifact{session.txt} files
are projections ending at research closeout. The HTML advertises
34,102 events and retains 4,439 event elements, including 3,616 token
notifications. Every one of those token observations matches the
corresponding JSONL line, including all cumulative and last-use fields.
The full record adds individual timestamps, tool observations,
configuration records and request-level usage. Its digest and the
reconciliation procedure are supplied in Appendix~\ref{app:reproduction}.

Within the research window, completed-item notifications record
4,350 command or polling observations, 676 file-change items, 357
image views, and 1,293 web-tool observations. The web count includes
searches and page retrieval; it is not a count of search queries.
There are also four completed sleep-tool observations. We count each
completed item once by its identifier, keeping these counts separate
from the outer code-tool calls that can contain several operations.
Of the command observations, 251 have nonzero exit status; these include
expected no-match results and interrupted processes, and cannot be
interpreted as 251 failed mathematical experiments.

The dated visible messages support the narrative below. Completed
human-message items record the GAP correction during research and the
model confirmation during manuscript preparation; the original
objectives also occur in goal records. These counts do not measure
oversight outside the recorded interaction. The derived public data
omit internal reasoning and system/developer instruction bodies.

\subsection{Token and cost accounting}

The JSONL contains two distinct usage ledgers. Its 3,829 unique
\texttt{token\_usage\_record} responses sum exactly to every running
\texttt{thread\_token\_usage} value. Of these responses, 84 match
context-compaction identifiers. Removing those responses gives exactly
the final cumulative \texttt{event\_msg/token\_count} counters, which
are the counter stream used by ccusage~\cite{ccusage-docs}.
Embedded copies of history and usage are not new requests.

The organizers supplied the ccusage row reproduced in
Table~\ref{tab:tokens}. It agrees exactly with the second ledger for
the \emph{whole supplied session}, including manuscript preparation.
Its input column counts uncached input; the renderer's input counter
includes cached input. Reasoning tokens are a subset of output, and
are not added a second time~\cite{promptcaching,reasoning}.

\begin{table}[ht]
\centering
\caption{Supplied ccusage row and its full-session reconciliation.
The row excludes the separately recorded compaction requests.}
\label{tab:tokens}
\begin{tabular}{lr}
\toprule
Counter & Tokens\\
\midrule
Uncached input & 21,605,689\\
Cached input & 482,598,400\\
Output & 4,055,235\\
\quad Of which reasoning output & 1,894,719\\
Total (uncached input plus cached input plus output) & 508,259,324\\
\midrule
Additional input and output in 84 compaction requests & 21,569,265\\
All recorded requests, including compactions & 529,828,589\\
\bottomrule
\end{tabular}
\end{table}

Using flat rates of USD 10 per million uncached input tokens, USD 1
per million cached input tokens, and USD 50 per million output tokens
reproduces the supplied USD 901.42 estimate:
\[
 \frac{10(21{,}605{,}689)+482{,}598{,}400+50(4{,}055{,}235)}{10^6}
 =901.41704.
\]
These rates match the model documentation consulted on 15 September
2026~\cite{astra}. The maximum recorded request input is 259,244
tokens, below its 272,000-token long-input threshold. We use the same
flat rates throughout Table~\ref{tab:phase-usage}; this is a transparent
API-equivalent reconstruction, not an invoice or an inference about
the organizers' subscription bill, discounts, or additional tool charges.
The exact ccusage version and invocation were not supplied.

\begin{table}[ht]
\centering\small
\caption{Usage by phase, assigned by request-record timestamp.
``Other'' means responses other than context compactions.
USD estimates are rounded independently to cents.}
\label{tab:phase-usage}
\begin{tabular}{lrrr}
\toprule
 & Research & Closeout & Manuscript\\
\midrule
Other responses & 3,515 & 5 & 225\\
Compaction responses & 79 & 0 & 5\\
All input and output tokens & 497,160,051 & 460,875 & 32,207,663\\
Other-response estimate (USD) & 846.09 & 0.56 & 54.77\\
Compaction estimate (USD) & 64.52 & 0.00 & 3.85\\
All-response estimate (USD) & 910.61 & 0.56 & 58.62\\
\bottomrule
\end{tabular}
\end{table}

Thus the full-session estimate including compactions is USD 969.78,
of which USD 910.61 is assigned to the 48-hour research window.
The earlier HTML total of 477,334,181 tokens equals research plus
closeout usage excluding compactions. A delayed token notification
straddles the deadline, so phase assignment uses request-record times,
not notification times. These are logging boundaries, not exact
provider execution intervals. Cached input dominates both ledgers;
neither total measures distinct text read or newly generated text.

\subsection{Outcome categories and checks}

We distinguish five questions about an artifact:
\begin{enumerate}
\item Does its mathematical statement match the question being addressed?
\item Is the argument correct under its stated assumptions and imported results?
\item Do the computations establish the particular finite claims assigned to them?
\item Is the assertion already known or an immediate consequence of prior work?
\item Has someone independent of the producing agent accepted the argument?
\end{enumerate}
These questions require different evidence. A digest match answers
none of the mathematical questions. A finite check cannot establish
an unrestricted infinite-group theorem without a proved reduction.
Failure to find a paper in a literature search does not establish
priority. A second implementation can expose an error in a first
program while sharing the same mistaken mathematical specification.

The deadline ledger uses \emph{complete candidate} for an argument
presented as answering the recorded scope and surviving the internal
checks conducted during the run. A \emph{partial result} proves a
restricted assertion while leaving the full question open.
A \emph{rediscovery} reproduces a known phenomenon or proof after
independent work, whereas a \emph{prior-result deduction} applies an
identified theorem. A bounded search without a hit remains
\emph{computational evidence} for its exact range.
These are reporting categories, not interchangeable levels on a
single scale of confidence.
